\section{Sterowanie zasilaniem}
Struktura \texttt{PowerManager} (przedstawiona wraz z~implementacją 
na listingu~\ref{lst:power_manager_struct}) odpowiada za sterowanie zasilaniem 
peryferiów o~podwyższonym poborze prądu (moduł RF, licznik 
Geigera, magistrala I2C), poprzez niezależne przełączanie 
odpowiadających im pinów cyfrowych skonfigurowanych jako 
wyjścia. Typy pinów są parametrami generycznymi, co pozwala 
uniknąć narzutu dynamicznego dyspozytora przy jednoczesnym 
zachowaniu elastyczności konfiguracji sprzętowej.

Istotnym elementem implementacji jest obsługa flagi \texttt{active\_low}, 
uwzględniającej fakt, że w~zależności od zastosowanego 
układu sterującego zasilaniem (np. tranzystora), stan 
aktywny modułu może odpowiadać zarówno stanowi wysokiemu, 
jak i~niskiemu na pinie sterującym. Dzięki zastosowaniu operacji 
XOR (\texttt{\^}) pomiędzy pożądanym stanem logicznym a~wartością tej 
flagi, interfejs udostępniany przez metody \texttt{activate\_*}/\texttt{deactivate\_* }
pozostaje niezależny od szczegółów elektrycznych konkretnego 
rozwiązania sprzętowego.

Funkcjonalność została podzielona na dwa poziomy granularności: 
metody \texttt{activate\_power\_hungry}/\texttt{deactivate\_power\_hungry} sterują 
wyłącznie modułem RF, którego pobór prądu jest na tyle duży, 
że wymaga on osobnej, częstszej kontroli niezależnie od 
pozostałych czujników, natomiast metody 
\texttt{activate\_all}/\texttt{deactivate\_all} zbiorczo przełączają wszystkie 
zarządzane piny za pośrednictwem prywatnej metody pomocniczej 
\texttt{set\_all\_state}, co pozwala na jednoczesne, całkowite odcięcie 
zasilania peryferiów.

\begin{lstlisting}[
    language=Rust, style=colouredRust, 
    caption={Struktura \texttt{Transmitter} wraz z~implementacją},
    label={lst:power_manager_struct}
]
pub struct PowerManager<GeigerPin, I2cPin, RfPin>{
    rf_vcc_pin: Pin<Output, RfPin>,
    geiger_vcc_pin: Pin<Output, GeigerPin>,
    i2c_vcc_pin: Pin<Output, I2cPin>,
    active_low: bool,
}

impl<GeigerPin, I2cPin, RfPin> PowerManager<GeigerPin, I2cPin, RfPin>
where GeigerPin: PinOps, I2cPin: PinOps, RfPin: PinOps{
    pub fn new(
        rf_vcc_pin: Pin<Output, RfPin>,
        geiger_vcc_pin: Pin<Output, GeigerPin>,
        i2c_vcc_pin: Pin<Output, I2cPin>,
        active_low: bool,
    ) -> Self{
        Self{rf_vcc_pin, geiger_vcc_pin, i2c_vcc_pin, active_low}
    }

    pub fn deactivate_power_hungry(&mut self){
        let state = (false ^ self.active_low).into();
        let _ = self.rf_vcc_pin.set_state(state);
    }

    pub fn activate_power_hungry(&mut self){
        let state = (true ^ self.active_low).into();
        let _ = self.rf_vcc_pin.set_state(state);
    }

    pub fn deactivate_all(&mut self){
        self.set_all_state((false ^ self.active_low).into());
    }

    pub fn activate_all(&mut self){
        self.set_all_state((true ^ self.active_low).into());
    }

    fn set_all_state(&mut self, state: PinState){
        let _ = self.rf_vcc_pin.set_state(state);
        let _ = self.geiger_vcc_pin.set_state(state);
        let _ = self.i2c_vcc_pin.set_state(state);
    }
}
\end{lstlisting}